%% LyX 2.4.4 created this file.  For more info, see https://www.lyx.org/.
%% Do not edit unless you really know what you are doing.
\documentclass[english]{article}
\usepackage[T1]{fontenc}
\usepackage[latin9]{inputenc}
\usepackage{enumitem}
\usepackage{amsmath}
\usepackage{amsthm}
\usepackage{amssymb}
\usepackage{geometry}
\geometry{verbose,tmargin=1in,bmargin=1in,lmargin=1in,rmargin=1in}

\makeatletter
%%%%%%%%%%%%%%%%%%%%%%%%%%%%%% Textclass specific LaTeX commands.
\numberwithin{equation}{section}
\numberwithin{figure}{section}
\newlength{\lyxlabelwidth}      % auxiliary length 
\theoremstyle{definition}
\newtheorem*{defn*}{\protect\definitionname}
\theoremstyle{plain}
\newtheorem*{lem*}{\protect\lemmaname}

%%%%%%%%%%%%%%%%%%%%%%%%%%%%%% User specified LaTeX commands.
\usepackage{fancyhdr}
\usepackage{lastpage}
\pagestyle{fancy}
\usepackage{enumitem}
\usepackage{ifthen}
\everymath=\expandafter{\the\everymath\displaystyle}
%\DeclareMathSizes{10}{10}{10}{10}
\usepackage{multicol}

\usepackage{tikz}
\usetikzlibrary{patterns}
\usetikzlibrary{plotmarks}
\usepackage{pgfplots}
\pgfplotsset{compat=newest}

\usepackage{calc}
\usepackage[nomessages]{fp}

\usepackage{datetime}
% Added by lyx2lyx
\setlength{\parskip}{\smallskipamount}
\setlength{\parindent}{0pt}

\makeatother

\usepackage{babel}
\providecommand{\definitionname}{Definition}
\providecommand{\lemmaname}{Lemma}

\begin{document}
\renewcommand{\author}{Dr.\,James Ashe}

\newcommand{\classNum}{335}

\newcommand{\classSec}{A}

\newcommand{\nameType}{Name: \hspace{-4cm}}

\newcommand{\examType}{Homework 5}

\newcommand{\Date}{\today} %%\AdvanceDate[12] \today

%\newdate{my_date}{\day}{\month}{\year}%   
%\AdvanceDate[-5] 
%\displaydate{my_date}

%%First page's header%%%%%%%%%%%%%%%%%%%%%%
\fancypagestyle{firststyle} %%header settings for first pagr
{
	\fancyhead[L]{MTH \classNum{}(\classSec{}) \author{}\\
	\textbf{\examType{}} \Date{}}
	\fancyhead[C]{\textbf{\nameType}}
	\setlength{\headsep}{14pt}
}
%%%%%%%%%%%%%%%%%%%%%%%%%%%%%%%%%%%%%%%%%%

%%Next page's header%%%%%%%%%%%%%%%%%%%%%%%
%\setlist{nolistsep} %eliminates space between list items
\lhead{MTH \classNum{}(\classSec{})} 
\chead{\textbf{\nameType}} 
\cfoot{\thepage{}~of~\pageref{LastPage}} 
\setlength{\headsep}{5pt} 
%%%%%%%%%%%%%%%%%%%%%%%%%%%%%%%%%%%%%%%%%%%

%%Various settings; see the preamble for other settings%%%%%
%%\setenumerate[0]{leftmargin=12pt,itemindent=0pt,labelwidth=0pt}
\renewcommand{\labelenumi}{\textbf{\arabic{enumi}.}}
\renewcommand{\labelenumii}{\textbf{\alph{enumii}.}}
\thispagestyle{firststyle}
%%%%%%%%%%%%%%%%%%%%%%%%%%%%%%%%%%%%%%%%%%

\textbf{Homework Problems. }\emph{Bonus Problems}
\begin{defn*}
A \emph{group homomorphism} is a function $f:G_{1}\rightarrow G_{2}$
between two groups that preserves the group operation; i.e., $f(ab)=\underset{\mbox{operation of }G_{2}}{\underbrace{f(a)f(b)}\mbox{.}}$
\end{defn*}
\begin{itemize}
\item [ex.]The mapping $f:\mathbb{Z}\rightarrow\mathbb{Z}_{n}$ where $\mathbb{Z}$
and $\mathbb{Z}_{n}$ are groups under addition and modular addition
respectively. Since $f(a+b)=(a+b)\bmod n=a\bmod n+b\bmod n$.
\end{itemize}
\begin{defn*}
Let $G$ be a group a $H$ a subset of $G$. For any $a\in G$, the
set $aH=\left\{ ah|h\in H\right\} $ is called the \emph{left coset
of H in G containing a. }$Ha=\left\{ ha|h\in H\right\} $ is called
the \emph{right coset of H in G containing a}. 
\end{defn*}
\begin{lem*}
For $H$ a subgroup of $G$ and $a,b\in G$, the following properties
of cosets hold:

\begin{enumerate}
\item $a\in aH$.\vfill
\item $aH=H$ if and only if $a\in H$. (parts 1 and 2 must be done together)
\item $aH=bH$ or $aH\cap bH=\emptyset$.
\item $\left|aH\right|=\left|Ha\right|=\left|bH\right|=\left|H\right|$.\vfill
\end{enumerate}
\end{lem*}
\begin{enumerate}[resume]
\item \textbf{Lagrange's Theorem: $\left|H\right||\left|G\right|$. }If
$G$ is a finite group and $H$ is a subgroup of $G$, then $\left|H\right|$
divides $\left|G\right|$. (also $\frac{\left|G\right|}{\left|H\right|}$is
the number of distinct left cosets in $G$).\vfill
\item Prove that if $\left|G\right|\leq4$ then $G$ is abelian. Then prove
that if $\left|G\right|<6$ then $G$ is abelian. \vfill
\item Let $f:G_{1}\rightarrow G_{2}$ be a group homomorphism and let $H\leq G_{2}$.
Prove that $f^{-1}(H)=\left\{ g\in G_{1}|f(g)\in H\right\} $ is a
subgroup of $G_{1}$.\vfill
\end{enumerate}

\end{document}
